% GENERATED FILE. Edit canonical lessons, then run npm run generate:books.
% canonical-source-hash: fnv1a64:3d5c08bbfb658680
% canonical-lessons: MR-C557-kuvat, MR-C557-avaka, MR-C557-katkasar, MR-C557-purvasucana, MR-C557-anumati, MR-R557-review-551-557, MR-R557-review-final-557

\chapter{Means and Notice}
\label{ch:mr-means-and-notice}
\hlchaptermodality{557}

\begin{chapteropening}
\textbf{By the end of this chapter:} \emph{I can talk about means, scope, thrift, notice and permission.}

Complete the last lesson of chapter 557: I can talk about means, scope, thrift, notice and permission.
\end{chapteropening}

\section[\texorpdfstring{\mr{कुवत} (kuvat)}{कुवत (kuvat)}]{\mr{कुवत} (kuvat) — capacity; means}

\label{lesson:MR-C557-kuvat}



\begin{quote}
Before the new one: say the Marathi for dire, then the Marathi for confusion.
\end{quote}

\subsection*{You'll want to know: \mr{कुवत}}
\textbf{\mr{कुवत}} — \emph{kuvat} — \textquotedblleft{}capacity; means\textquotedblright{}.

\begin{quote}
\mr{घर} \mr{घेण्याची} \mr{आमची} \mr{अजून} \mr{कुवत} \mr{नाही}.
\end{quote}

\begin{quote}
\emph{ghar gheṇyācī āmcī ajūn kuvat nāhī.}
\end{quote}

\begin{quote}
We do not yet have the means to buy a house.
\end{quote}

Feminine.

\begin{cousinweb}
\textbf{\mr{कुवत}} (\emph{kuvat}) is a loan through Persian usually traced to Arabic \emph{quwwa}, \textquotedblleft{}strength\textquotedblright{}.
\end{cousinweb}

\begin{grammarlens}[title={What you've built}]
One more word for saying what happens if.
\end{grammarlens}

\subsection*{Your turn}
\begin{itemize}
\raggedright
  \item \emph{Say it:} \emph{kuvat}
  \item \emph{Say it:} \emph{kuvat}, once more
  \item \emph{Say it:} \emph{ghar gheṇyācī āmcī ajūn kuvat nāhī.} — We do not yet have the means to buy a house.
  \item \emph{Recall:} say the Marathi for dire, then the Marathi for confusion, then say \emph{kuvat} again
\end{itemize}

\begin{tcolorbox}[breakable,colback=teal!4,colframe=teal!35!black,title={Before you move on}]
What does \mr{कुवत} mean? (\textquotedblleft{}Capacity; means\textquotedblright{}.) Say it once more.
\end{tcolorbox}
\section[\texorpdfstring{\mr{आवाका} (āvākā)}{आवाका (āvākā)}]{\mr{आवाका} (āvākā) — scope; reach}

\label{lesson:MR-C557-avaka}



\begin{quote}
Before the new one: say the Marathi for confusion, then the Marathi for capacity.
\end{quote}

\subsection*{You'll want to know: \mr{आवाका}}
\textbf{\mr{आवाका}} — \emph{āvākā} — \textquotedblleft{}scope; reach\textquotedblright{}.

\begin{quote}
\mr{हे} \mr{काम} \mr{माझ्या} \mr{आवाक्याबाहेर} \mr{आहे}.
\end{quote}

\begin{quote}
\emph{he kām mājhyā āvākyābāher āhe.}
\end{quote}

\begin{quote}
This job is beyond my reach.
\end{quote}

Masculine.

The oblique is \mr{आवाक्या}-: \mr{आवाक्याबाहेर}, beyond reach.

\begin{cousinweb}
\textbf{\mr{आवाका}} (\emph{āvākā}) is a Marathi word for the reach of one's ability; its origin is not settled.
\end{cousinweb}

\begin{grammarlens}[title={What you've built}]
One more word for saying what happens if.
\end{grammarlens}

\subsection*{Your turn}
\begin{itemize}
\raggedright
  \item \emph{Say it:} \emph{āvākā}
  \item \emph{Say it:} \emph{āvākā}, once more
  \item \emph{Say it:} \emph{he kām mājhyā āvākyābāher āhe.} — This job is beyond my reach.
  \item \emph{Recall:} say the Marathi for confusion, then the Marathi for capacity, then say \emph{āvākā} again
\end{itemize}

\begin{tcolorbox}[breakable,colback=teal!4,colframe=teal!35!black,title={Before you move on}]
What does \mr{आवाका} mean? (\textquotedblleft{}Scope; reach\textquotedblright{}.) Say it once more.
\end{tcolorbox}
\section[\texorpdfstring{\mr{काटकसर} (kāṭkasar)}{काटकसर (kāṭkasar)}]{\mr{काटकसर} (kāṭkasar) — thrift}

\label{lesson:MR-C557-katkasar}



\begin{quote}
Before the new one: say the Marathi for capacity, then the Marathi for scope.
\end{quote}

\subsection*{You'll want to know: \mr{काटकसर}}
\textbf{\mr{काटकसर}} — \emph{kāṭkasar} — \textquotedblleft{}thrift\textquotedblright{}.

\begin{quote}
\mr{आजीने} \mr{काटकसर} \mr{करून} \mr{थोडे} \mr{पैसे} \mr{वाचवले}.
\end{quote}

\begin{quote}
\emph{ājīne kāṭkasar karūn thoḍe paise vācavle.}
\end{quote}

\begin{quote}
By being thrifty Grandmother saved a little money.
\end{quote}

Feminine.

\mr{काटकसर} \mr{करणे}, to economize.

\begin{cousinweb}
\textbf{\mr{काटकसर}} (\emph{kāṭkasar}) is a Marathi word for careful spending; its origin is not settled.
\end{cousinweb}

\begin{grammarlens}[title={What you've built}]
One more word for saying what happens if.
\end{grammarlens}

\subsection*{Your turn}
\begin{itemize}
\raggedright
  \item \emph{Say it:} \emph{kāṭkasar}
  \item \emph{Say it:} \emph{kāṭkasar}, once more
  \item \emph{Say it:} \emph{ājīne kāṭkasar karūn thoḍe paise vācavle.} — By being thrifty Grandmother saved a little money.
  \item \emph{Recall:} say the Marathi for capacity, then the Marathi for scope, then say \emph{kāṭkasar} again
\end{itemize}

\begin{tcolorbox}[breakable,colback=teal!4,colframe=teal!35!black,title={Before you move on}]
What does \mr{काटकसर} mean? (\textquotedblleft{}Thrift\textquotedblright{}.) Say it once more.
\end{tcolorbox}
\section[\texorpdfstring{\mr{पूर्वसूचना} (pūrvasūcanā)}{पूर्वसूचना (pūrvasūcanā)}]{\mr{पूर्वसूचना} (pūrvasūcanā) — advance notice}

\label{lesson:MR-C557-purvasucana}



\begin{quote}
Before the new one: say the Marathi for scope, then the Marathi for thrift.
\end{quote}

\subsection*{You'll want to know: \mr{पूर्वसूचना}}
\textbf{\mr{पूर्वसूचना}} — \emph{pūrvasūcanā} — \textquotedblleft{}advance notice\textquotedblright{}.

\begin{quote}
\mr{कोणतीही} \mr{पूर्वसूचना} \mr{न} \mr{देता} \mr{वीज} \mr{गेली}.
\end{quote}

\begin{quote}
\emph{koṇtīhī pūrvasūcanā na detā vīj gelī.}
\end{quote}

\begin{quote}
The power went off without any advance notice.
\end{quote}

Feminine.

\begin{cousinweb}
\textbf{\mr{पूर्वसूचना}} (\emph{pūrvasūcanā}) is a tatsama compound of Sanskrit \emph{pūrva}, \textquotedblleft{}before\textquotedblright{}, and \emph{sūcanā}, \textquotedblleft{}indication\textquotedblright{} (Monier-Williams).
\end{cousinweb}

\begin{grammarlens}[title={What you've built}]
One more word for saying what happens if.
\end{grammarlens}

\subsection*{Your turn}
\begin{itemize}
\raggedright
  \item \emph{Say it:} \emph{pūrvasūcanā}
  \item \emph{Say it:} \emph{pūrvasūcanā}, once more
  \item \emph{Say it:} \emph{koṇtīhī pūrvasūcanā na detā vīj gelī.} — The power went off without any advance notice.
  \item \emph{Recall:} say the Marathi for scope, then the Marathi for thrift, then say \emph{pūrvasūcanā} again
\end{itemize}

\begin{tcolorbox}[breakable,colback=teal!4,colframe=teal!35!black,title={Before you move on}]
What does \mr{पूर्वसूचना} mean? (\textquotedblleft{}Advance notice\textquotedblright{}.) Say it once more.
\end{tcolorbox}
\section[\texorpdfstring{\mr{अनुमती} (anumatī)}{अनुमती (anumatī)}]{\mr{अनुमती} (anumatī) — formal permission}

\label{lesson:MR-C557-anumati}



\begin{quote}
Before the new one: say the Marathi for thrift, then the Marathi for advance notice.
\end{quote}

\subsection*{You'll want to know: \mr{अनुमती}}
\textbf{\mr{अनुमती}} — \emph{anumatī} — \textquotedblleft{}formal permission\textquotedblright{}.

\begin{quote}
\mr{शिक्षकांची} \mr{अनुमती} \mr{घेऊन} \mr{मी} \mr{वर्गाबाहेर} \mr{गेलो}.
\end{quote}

\begin{quote}
\emph{śikṣakāñcī anumatī gheūn mī vargābāher gelo.}
\end{quote}

\begin{quote}
With the teacher's permission I went out of the class.
\end{quote}

Feminine.

\mr{अनुमती} is the formal word beside \mr{परवानगी}.

\begin{cousinweb}
\textbf{\mr{अनुमती}} (\emph{anumatī}) is a tatsama, taken from Sanskrit \emph{anumati}, \textquotedblleft{}assent, permission\textquotedblright{} (Monier-Williams).
\end{cousinweb}

\begin{grammarlens}[title={What you've built}]
One more word for saying what happens if.
\end{grammarlens}

\subsection*{Your turn}
\begin{itemize}
\raggedright
  \item \emph{Say it:} \emph{anumatī}
  \item \emph{Say it:} \emph{anumatī}, once more
  \item \emph{Say it:} \emph{śikṣakāñcī anumatī gheūn mī vargābāher gelo.} — With the teacher's permission I went out of the class.
  \item \emph{Recall:} say the Marathi for thrift, then the Marathi for advance notice, then say \emph{anumatī} again
\end{itemize}

\begin{tcolorbox}[breakable,colback=teal!4,colframe=teal!35!black,title={Before you move on}]
What does \mr{अनुमती} mean? (\textquotedblleft{}Formal permission\textquotedblright{}.) Say it once more.
\end{tcolorbox}
\section[(dialogue)]{Looking Back at Chances and Choices}

\label{lesson:MR-R557-review-551-557}



\begin{quote}
Say the Marathi for waiting and a list.
\end{quote}

\subsection*{Your turn}
Say these aloud:

\begin{itemize}
\raggedright
  \item waiting, a list, a sequence, reserved, unforeseen
  \item to fall apart, to fizzle out, to crumble, to break down, to come loose
  \item to achieve, to elude, accurate, without fail, prompt
  \item on purpose, in advance, at the last moment, mostly, roughly
  \item a span of time, a needless delay, urgency, a rush, a scramble
  \item favourable, unfavourable, a muddle, dire, confusion
  \item capacity, scope, thrift, advance notice, formal permission
\end{itemize}

\begin{tcolorbox}[breakable,colback=teal!4,colframe=teal!35!black,title={Before you move on}]
Waiting? (\textbf{\mr{प्रतीक्षा}}, \emph{pratīkṣā}.) A list? (\textbf{\mr{यादी}}, \emph{yādī}.) A sequence? (\textbf{\mr{क्रम}}, \emph{kram}.) Reserved? (\textbf{\mr{राखीव}}, \emph{rākhīv}.) Unforeseen? (\textbf{\mr{आकस्मिक}}, \emph{ākasmik}.) To fall apart? (\textbf{\mr{फिसकटणे}}, \emph{phiskaṭṇe}.) To fizzle out? (\textbf{\mr{बारगळणे}}, \emph{bārgaḷṇe}.) To crumble? (\textbf{\mr{ढासळणे}}, \emph{ḍhāsaḷṇe}.) To break down? (\textbf{\mr{कोलमडणे}}, \emph{kolmaḍṇe}.) To come loose? (\textbf{\mr{सुटणे}}, \emph{suṭṇe}.) To achieve? (\textbf{\mr{साध्य} \mr{करणे}}, \emph{sādhya karṇe}.) To elude? (\textbf{\mr{हुलकावणी} \mr{देणे}}, \emph{hulkāvṇī deṇe}.) Accurate? (\textbf{\mr{अचूक}}, \emph{acūk}.) Without fail? (\textbf{\mr{हमखास}}, \emph{hamkhās}.) Prompt? (\textbf{\mr{तत्पर}}, \emph{tatpar}.) On purpose? (\textbf{\mr{मुद्दाम}}, \emph{muddām}.) In advance? (\textbf{\mr{आगाऊ}}, \emph{āgāū}.) At the last moment? (\textbf{\mr{ऐनवेळी}}, \emph{ainveḷī}.) Mostly? (\textbf{\mr{बहुतेक}}, \emph{bahutek}.) Roughly? (\textbf{\mr{साधारण}}, \emph{sādhāraṇ}.) A span of time? (\textbf{\mr{अवधी}}, \emph{avadhī}.) A needless delay? (\textbf{\mr{दिरंगाई}}, \emph{diraṅgāī}.) Urgency? (\textbf{\mr{तातडी}}, \emph{tātḍī}.) A rush? (\textbf{\mr{धांदल}}, \emph{dhāndal}.) A scramble? (\textbf{\mr{तारांबळ}}, \emph{tārāmbaḷ}.) Favourable? (\textbf{\mr{अनुकूल}}, \emph{anukūl}.) Unfavourable? (\textbf{\mr{प्रतिकूल}}, \emph{pratikūl}.) A muddle? (\textbf{\mr{घोळ}}, \emph{ghoḷ}.) Dire? (\textbf{\mr{बिकट}}, \emph{bikaṭ}.) Confusion? (\textbf{\mr{संभ्रम}}, \emph{sambhram}.) Capacity? (\textbf{\mr{कुवत}}, \emph{kuvat}.) Scope? (\textbf{\mr{आवाका}}, \emph{āvākā}.) Thrift? (\textbf{\mr{काटकसर}}, \emph{kāṭkasar}.) Advance notice? (\textbf{\mr{पूर्वसूचना}}, \emph{pūrvasūcanā}.) Formal permission? (\textbf{\mr{अनुमती}}, \emph{anumatī}.)
\end{tcolorbox}
\section[(dialogue)]{Means and Notice — once more}

\label{lesson:MR-R557-review-final-557}



\begin{quote}
Say the Marathi for capacity and scope.
\end{quote}

\subsection*{Your turn}
\begin{itemize}
\raggedright
  \item \emph{Say it:} capacity, scope, thrift, advance notice, formal permission
\end{itemize}

\begin{tcolorbox}[breakable,colback=teal!4,colframe=teal!35!black,title={Before you move on}]
Capacity? (\textbf{\mr{कुवत}}, \emph{kuvat}.) Scope? (\textbf{\mr{आवाका}}, \emph{āvākā}.) Thrift? (\textbf{\mr{काटकसर}}, \emph{kāṭkasar}.) Advance notice? (\textbf{\mr{पूर्वसूचना}}, \emph{pūrvasūcanā}.) Formal permission? (\textbf{\mr{अनुमती}}, \emph{anumatī}.)
\end{tcolorbox}
